% Alternatives for the elements of figure 12 (momentum representation). Three approaches per element.
\documentclass[10pt,border=2mm]{standalone}
\usepackage[T1]{fontenc}
\usepackage{figurestyle}
\input{primitives.tex}
\input{numbers.tex}
\begin{document}
\begin{tikzpicture}[plate]
\path[use as bounding box] (0,0) rectangle (15,24);
% ================================================= element 1: the K ladder
\node[lab,anchor=west] at (.2,23.6) {Element 1: the $K$ ladder (body-axis component of $\mathbf J$)};
% (i) nested cones, Harter 5.4.4
\begin{scope}[shift={(2.6,20.6)}]
  \pgfmathsetmacro{\s}{.55}\pgfmathsetmacro{\E}{.32}
  \draw[fine] (0,{-2.6*\s})--(0,{2.9*\s});
  \foreach \K in {2,1,0,-1,-2} {\pgfmathsetmacro{\r}{sqrt(6-\K*\K)*\s}\pgfmathsetmacro{\h}{\K*\s}
    \draw[fine] (-\r,\h) -- (0,0) -- (\r,\h); \draw[fine] (0,\h) ellipse[x radius=\r,y radius={\E*\r}];
    \node[sub,anchor=east] at ({-\r-.15},\h) {$\K$};}
  \draw[harrow] (0,0)--({sqrt(5)*\s*cos(20)},{\s+\E*sqrt(5)*\s*sin(20)});
  \node[sub,anchor=north,text width=4cm,align=center] at (0,-1.7) {(i) nested cones, after Harter 5.4.4: every $K$ at once, the cone family};
\end{scope}
% (ii) the sphere of radius |J| with latitude rings at heights K, Harter 5.5.3
\begin{scope}[shift={(7.5,20.6)}]
  \pgfmathsetmacro{\s}{.55}\pgfmathsetmacro{\R}{sqrt(6)*\s}\pgfmathsetmacro{\E}{.32}
  \hatom{0}{0}{\R}{white}
  \foreach \K in {2,1,0,-1,-2} {\pgfmathsetmacro{\r}{sqrt(6-\K*\K)*\s}\pgfmathsetmacro{\h}{\K*\s}
    \draw[fine] ({-\r},\h) arc[start angle=180,end angle=360,x radius=\r,y radius={\E*\r}];
    \draw[hidden] ({-\r},\h) arc[start angle=180,end angle=0,x radius=\r,y radius={\E*\r}];
    \node[sub,anchor=west] at ({\R+.1},\h) {$K=\K$};}
  \draw[fine] (0,{-\R-.2})--(0,{\R+.3});
  \draw[harrow] (0,0)--({sqrt(5)*\s*cos(30)},{\s+\E*sqrt(5)*\s*sin(30)});
  \node[sub,anchor=north,text width=4cm,align=center] at (0,-1.7) {(ii) the sphere $|\mathbf J|=\sqrt{J(J+1)}$ with latitude rings at height $K$, after 5.5.3};
\end{scope}
% (iii) the minimal ladder: rungs and one tilted arrow
\begin{scope}[shift={(12.3,20.6)}]
  \pgfmathsetmacro{\s}{.55}
  \draw[fine] (0,{-2.6*\s})--(0,{2.9*\s});
  \foreach \K in {2,1,0,-1,-2} {\draw[heavy] (-.35,{\K*\s})--(.35,{\K*\s}); \node[sub,anchor=west] at (.45,{\K*\s}) {$K=\K$};}
  \draw[harrow] (0,0)--({sqrt(5)*\s*.8},{\s});
  \draw[fine,dash pattern=on .6pt off 1.3pt] ({sqrt(5)*\s*.8},{\s})--(0,{\s});
  \node[sub,anchor=north,text width=4cm,align=center] at (0,-1.7) {(iii) the ladder alone: rungs at $K$, one $\mathbf J$ projected; least ink, no cone family};
\end{scope}
% ================================================= element 2: the seam shift kappa = m + rho K
\node[lab,anchor=west] at (.2,17.4) {Element 2: the seam shift $\kappa=m+\rho K$ (twisted frame against periodic frame)};
% (i) dot rows with a gold shear line (current)
\begin{scope}[shift={(.6,14.6)}]
  \foreach \K in {-2,...,2} {\pgfmathsetmacro{\y}{1.2+.42*\K}\node[sub,anchor=east] at (-.1,\y) {$\K$}; \draw[fine] (0,\y)--(3.3,\y);
    \foreach \m in {-3,...,3} \fill[PlateInk] ({1.65+.42*(\m+\rhoIll*\K)},\y) circle[radius=1.3pt];}
  \draw[line width=1.1pt,draw=ColGold] ({1.65-.42*\rhoIll*2.4},{1.2-.42*2.4})--({1.65+.42*\rhoIll*2.4},{1.2+.42*2.4});
  \node[sub,anchor=north,text width=4.2cm,align=center] at (1.65,-.1) {(i) dot rows, one per $K$, the $m=0$ dots joined by the gold line $\kappa=\rho K$};
\end{scope}
% (ii) a sheared grid: the periodic lattice of (m, K) as a rectangle, the twisted one as the sheared parallelogram
\begin{scope}[shift={(5.5,14.6)}]
  \path[fill=PlateInk!6] (0,.36)--(3.3,.36)--(3.3,2.04)--(0,2.04)--cycle;
  \draw[fine] (0,.36) rectangle (3.3,2.04);
  \draw[heavy] ({-.42*\rhoIll*2},.36)--({3.3-.42*\rhoIll*2},.36)--({3.3+.42*\rhoIll*2},2.04)--({.42*\rhoIll*2},2.04)--cycle;
  \foreach \K in {-2,...,2} {\pgfmathsetmacro{\y}{1.2+.42*\K}\foreach \m in {-3,...,3} {\fill[PlateInk!35] ({1.65+.42*\m},\y) circle[radius=1.1pt]; \fill[PlateInk] ({1.65+.42*(\m+\rhoIll*\K)},\y) circle[radius=1.3pt];}}
  \node[sub,anchor=north,text width=4.2cm,align=center] at (1.65,-.1) {(ii) the $(m,K)$ lattice as a sheared grid: grey the periodic frame, ink the twisted frame, outline sheared by $\rho$};
\end{scope}
% (iii) a single row, with the shift as a bracket
\begin{scope}[shift={(10.6,14.6)}]
  \foreach \K/\y in {1/1.7,0/1.2} {\node[sub,anchor=east] at (-.1,\y) {$K=\K$}; \draw[fine] (0,\y)--(3.3,\y);
    \foreach \m in {-3,...,3} {\fill[PlateInk!35] ({1.65+.42*\m},\y) circle[radius=1.1pt]; \fill[PlateInk] ({1.65+.42*(\m+\rhoIll*\K)},\y) circle[radius=1.3pt];}}
  \draw[fine] (1.65,1.95)--(1.65,2.1)--({1.65+.42*\rhoIll},2.1)--({1.65+.42*\rhoIll},1.95); \node[sub,anchor=south] at ({1.65+.21*\rhoIll},2.12) {$\rho K$};
  \node[sub,anchor=north,text width=4.2cm,align=center] at (1.65,-.1) {(iii) two rows only, the shift of one row bracketed; the rule, not the picture of all rows};
\end{scope}
% ================================================= element 3: the torsion species
\node[lab,anchor=west] at (.2,11.8) {Element 3: periodic torsion modes and their species under $G_6$};
% (i) waveforms in a strip with level rules (current)
\begin{scope}[shift={(.4,9.2)}]
  \foreach \m/\cs/\ss in \torsionSpecies {\pgfmathsetmacro{\cx}{.7+1.0*\m}
    \draw[line width=.6pt,draw=PlateInk] plot[domain=0:1,samples=50,variable=\x] ({\cx-.42+.84*\x},{1.5+.2*cos(360*\m*\x)});
    \draw[line width=.6pt,draw=PlateInk] plot[domain=0:1,samples=50,variable=\x] ({\cx-.42+.84*\x},{.9+.2*sin(360*\m*\x)});
    \node[sub] at (\cx,2.0) {$\m$};
    \ifnum\m=0 \draw[heavy] (\cx-.35,.45)--(\cx+.35,.45); \node[sub,anchor=north] at (\cx,.42) {$\mathrm A_1$};
    \else\ifnum\numexpr\m-3*(\m/3)\relax=0 \draw[heavy] (\cx-.35,.45)--(\cx+.35,.45); \node[sub,anchor=north] at (\cx,.42) {$\mathrm A_1$}; \draw[heavy] (\cx-.35,.05)--(\cx+.35,.05); \node[sub,anchor=north] at (\cx,.02) {$\mathrm A_2$};
    \else \draw[heavy] (\cx-.35,.48)--(\cx+.35,.48) (\cx-.35,.42)--(\cx+.35,.42); \node[sub,anchor=north] at (\cx,.39) {$\mathrm E$}; \fi\fi}
  \node[sub,anchor=north,text width=6.6cm,align=center] at (3.7,-.5) {(i) a strip of waveforms with level rules beneath, after 4.2.3};
\end{scope}
% (ii) the modes drawn on the torsion circle, so the three-fold symmetry is visible
\begin{scope}[shift={(8.6,9.9)}]
  \foreach \m/\dx in {1/0,2/1.7,3/3.4,6/5.1} {
    \begin{scope}[shift={(\dx,0)}]
      \draw[fine] (0,0) circle[radius=.5];
      \draw[line width=.7pt,draw=PlateInk] plot[domain=0:360,samples=120,variable=\a] ({(.5+.22*cos(\m*\a))*cos(\a)},{(.5+.22*cos(\m*\a))*sin(\a)});
      \foreach \a in {90,210,330} \draw[fine] ({.5*cos(\a)},{.5*sin(\a)})--({.62*cos(\a)},{.62*sin(\a)});
      \node[sub] at (0,-1.0) {$\cos\,\m\tau$};
    \end{scope}}
  \node[sub,anchor=north,text width=6.6cm,align=center] at (2.55,-1.2) {(ii) the mode as a radial plot on the torsion circle, with the three version marks: $3\mid m$ is visibly $t$-invariant ($\mathrm A_1$); $m=1,2$ are not};
\end{scope}
% (iii) a correlation diagram: m-levels on the left branching to species, after 4.2.3 literally
\begin{scope}[shift={(.6,3.2)}]
  \node[sub,anchor=south] at (.6,3.9) {$m$}; \node[sub,anchor=south] at (3.4,3.9) {species of $G_6$};
  \foreach \m/\cs/\ss in \torsionSpecies {\pgfmathsetmacro{\y}{.3+.55*\m}
    \draw[heavy] (.1,\y)--(1.1,\y); \node[sub,anchor=east] at (0,\y) {$\m$};
    \ifnum\m=0 \draw[fine,dash pattern=on 2pt off 1.2pt] (1.1,\y)--(2.9,\y); \draw[heavy] (2.9,\y)--(3.9,\y); \node[sub,anchor=west] at (4.0,\y) {$\mathrm A_1$};
    \else\ifnum\numexpr\m-3*(\m/3)\relax=0 \draw[fine,dash pattern=on 2pt off 1.2pt] (1.1,\y)--(2.9,{\y+.12}) (1.1,\y)--(2.9,{\y-.12}); \draw[heavy] (2.9,{\y+.12})--(3.9,{\y+.12}) (2.9,{\y-.12})--(3.9,{\y-.12}); \node[sub,anchor=west] at (4.0,{\y+.12}) {$\mathrm A_1$}; \node[sub,anchor=west] at (4.0,{\y-.12}) {$\mathrm A_2$};
    \else \draw[fine,dash pattern=on 2pt off 1.2pt] (1.1,\y)--(2.9,\y); \draw[heavy] (2.9,{\y+.03})--(3.9,{\y+.03}) (2.9,{\y-.03})--(3.9,{\y-.03}); \node[sub,anchor=west] at (4.0,\y) {$\mathrm E$}; \fi\fi}
  \node[sub,anchor=north,text width=5.6cm,align=center] at (2.4,-.2) {(iii) a correlation diagram in the manner of 4.2.3: levels $|m|$ at left, their species at right, no waveforms};
\end{scope}
\node[sub,anchor=north west,text width=7.4cm] at (7.4,6.8) {Critique written after rendering, see docs/illustration-notes.md};
\end{tikzpicture}
\end{document}
